\section{Introduction}
% The very first letter is a 2 line initial drop letter followed
% by the rest of the first word in caps.
% 
% form to use if the first word consists of a single letter:
% \IEEEPARstart{A}{demo} file is ....
% 
% form to use if you need the single drop letter followed by
% normal text (unknown if ever used by the IEEE):
% \IEEEPARstart{A}{}demo file is ....
% 
% Some journals put the first two words in caps:
% \IEEEPARstart{T}{his demo} file is ....
% 
% Here we have the typical use of a "T" for an initial drop letter
% and "HIS" in caps to complete the first word.
%\IEEEPARstart{A}{s} 
As Artificial Intelligence (AI) permeates our lives more and more, mechanisms such as deep neural networks~\cite{sze2017dnn} are put to use (or their deployment planned) in more and more places in various distributed systems. 

The data pipeline with AI typically requires the creation and the use of a \emph{model}, i.e. a layered structure of complex algorithms (also known as \emph{operators}) which interpret data and make decisions based on that data. This model must first be trained, before it can be put to use in production~\cite{sze2017dnn}.
Designing a distributed system combining AI and Internet of Things (AIoT) thus requires placing somewhere along the terminal-edge-cloud continuum both model training (learning) and model execution (inference).

Recent work on TinyML~\cite{sanchez2020tinyml,ray2022tinyml-review} forays into optimizing models to fit tinier resource budgets (and to perform efficiently nevertheless) on low-power microcontrollers. As a consequence, both learning placement possibilities, and inference placement possibilities are extended to encompass ultra low-power terminals.

However, generic and convenient open source tools lack for AIoT designers, who are required to:
\begin{itemize}
    \item evaluate the performance of their models when placed somewhere along the terminal-edge-cloud continuum, especially when including potential placement on different microcontroller-based devices;
    \item fine-tune their models, and identifying performance bottlenecks at model layer granularity, on different microcontrollers;
    \item select an adequate microcontroller to execute their model, for a targeted task running on a low-power device to-be-designed.
\end{itemize}

This paper thus introduces U-TOE, an open source AIoT toolkit which tightly combines a generic model compiler and a popular low-power IoT operating system to automatically compress, flash and evaluate arbitrary models (output of TensorFlow, PyTorch...) on arbitrary COTS low-power boards (such as BBC:microbit, nrf52840dk, Arduino Zero, HiFive...) based on the popular microcontroller architectures (ARM Cortex-M, RISC-V). 


{\bf Paper contributions.} Our contributions are given as follows:
\begin{itemize}
    % Some parts of them may should move to conclusion
    \item We designed a universal TinyML on-board evaluation toolkit for AIoT (U-TOE). It provides feasibility check of the deployment of arbitrary model on a given IoT hardware platform, allows researchers and developers to locate the performance bottleneck of a given model on target device. The evaluation results enable co-design with other components on system level, help optimizing ML models and configurations for specific use cases, allowing to achieve the best possible performance on target devices.
    
    \item We released the source code\footnote{see \url{https://github.com/zhaolanhuang/u-toe}} of U-TOE under an open licence. This implementation enables compilation, flashing, and evaluation of neural network (computational graph based) models from mainstream ML framework onto various low-power boards based on popular ISA.

    \item We provided benchmarks and a comparative experimental evaluation using U-TOE, reproducible on an open-access IoT testbed and on personal workstation, which both provides insights on inference performance with different models on different low-power hardware and
    demonstrates how U-TOE can be re-used by TinyML experimental researchers and developers to fine-tune AIoT configurations. %  the allows to compare the evaluation results on different hardware, and makes it easier for researchers and developers to decide the appropriate configurations.
\end{itemize}